% !TeX root = ../main.tex
%%% ---------------------------------------------------------------------------
%%% 摘要与关键词
%%%
%%% GB/T 7713.2 对摘要的要求（写之前先看一遍）：
%%%   · 报道性摘要，能脱离全文独立存在
%%%   · 不出现「本文」「作者」等自指表述
%%%   · 不引用文献编号、不出现图表号
%%%   · 不使用非公知的缩略语
%%%   · 中文摘要一般 300 字以内，英文摘要与中文对应
%%%   结构建议：目的 → 方法 → 结果（带具体数据）→ 结论
%%%
%%% 关键词写在 setup.tex 里，这里的环境会自动排出来。
%%% ---------------------------------------------------------------------------

\begin{abstract}
针对遥感图像中目标尺度跨度大、小目标特征在深层网络中易丢失的问题，提出一种
多尺度特征融合的目标检测方法。该方法在特征金字塔的基础上引入跨层注意力
门控，按空间位置自适应地加权来自不同下采样倍率的特征，并以可变形卷积
对齐相邻层级的感受野偏移。在 \dataset{} 数据集上的实验表明，所提方法的
平均精度均值达到 \qty{78.4}{\percent}，较基线方法提升 \qty{3.7}{\percent}，
其中小目标类别的平均精度提升 \qty{6.2}{\percent}，推理速度保持在
\qty{31}{\fps}。结果说明跨层注意力门控能够在不显著增加计算量的
前提下缓解小目标特征衰减，为高分辨率遥感图像的实时检测提供了可行方案。
\end{abstract}

%%% 英文摘要块（lang=en 时删掉这一段即可）
\begin{enabstract}
To address the large scale variation of objects and the loss of small-object
features in deep networks for remote sensing imagery, a multi-scale feature
fusion detection method is proposed. Built upon a feature pyramid, the method
introduces a cross-level attention gate that adaptively weights features from
different downsampling rates at each spatial location, and employs deformable
convolution to align the receptive-field offsets between adjacent levels.
Experiments on \dataset{} show that the proposed method attains a mean average
precision of \qty{78.4}{\percent}, an improvement of \qty{3.7}{\percent} over
the baseline, with a \qty{6.2}{\percent} gain on small-object categories while
maintaining an inference speed of \qty{31}{\fps}. The results
indicate that the cross-level attention gate mitigates small-object feature
attenuation without a significant increase in computation, offering a practical
route to real-time detection in high-resolution remote sensing imagery.
\end{enabstract}
